\section{Related Work}
\label{sec:related}

\paragraph{Skills and compact context representations.}
Skill benchmarks ask whether agents retrieve and apply reusable procedures
\citep{su2026srabench, skillsbench2026, sweskills2026}; we ask how a model
represents one already in context. Task and function vectors transfer behaviours
from in-context demonstrations through compact activations
\citep{hendel2023taskvectors, todd2024functionvectors}, and task arithmetic
makes related edits to weights \citep{ilharco2023taskarithmetic}.
\citet{dong2026taskvectorlimits} bound the task mappings a single vector can
express, including high-rank ones; our rank is instead that of a matrix of
hidden-state differences across skill positions, tested under several answer
formats. Others learn to compress context---aggregated activations
\citep{ardoin2026promptcompression}, summary tokens \citep{mu2023gist}, compact
states \citep{ge2024icae}, weight updates \citep{doc2lora2026,
deepcontextdistill2025}---while we examine the original model's states without
training a compression mechanism.

\paragraph{Answer representations and activation interventions.}
Multiple-choice answering requires selecting an answer symbol
\citep{wiegreffe2025mcq}, whereas reasoning can depend on states distributed
over generated positions \citep{dura2026sequentialpatching}. Transferring states
from completed reasoning \citep{mehrafarin2026cothidden} therefore tests a
different setting from ours, which captures states before generation begins.
Repeatedly applying activation directions can also steer extended generations
\citep{panickssery2024caa, liu2024icv, stolfo2025instruction}; our
single-position tests apply a one-time intervention at a prompt position.
Following activation-patching methodology \citep{zhang2023bestpractices,
heimersheim2024patching}, we compare generated answers against explicit
baselines and verify that writing a run's own states back leaves it unchanged.
Prior work traces factual information and context use across layers
\citep{meng2022rome, geva2023dissecting, sia2024incontexttranslation,
onset2025instruction}. We ask whether states stop transferring a skill's effect
at the same depth that the model stops depending on attention to them.

\paragraph{Representation rank and spectral measurements.}
Spectral analyses use relative singular-value sizes to study hidden-state
geometry \citep{valeriani2023geometry}, sensitivity to perturbations
\citep{sarkar2026causaldim}, and transformations across layers
\citep{fernando2026residualspectral}; we instead truncate the matrix of
content-specific state differences across skill positions and measure the
resulting performance. Reports of unusually large activations
\citep{sun2024massive, shi2026melayer} and dominant hidden dimensions
\citep{timkey2021rogue, kovaleva2021bertbusters} motivate a separate check:
does an apparent drop in dimensionality reflect a change in useful content or
merely a few large coordinates? We keep that geometric diagnostic distinct from
the rank intervention's behavioural outcome.

